\par\Needspace{6\baselineskip}\medskip
\begin{center}\hypertarget{native-236-article171}{}\pdfbookmark[2]{171.}{article-171}171.\end{center}
Zur Veranschaulichung dieser Untersuchung bietet sich ein Beispiel anhand der \emph{Vesta} an, die unter allen zuletzt entdeckten Planeten die geringste Neigung zur \SourceLeaf{237}Ekliptik aufweist\footnote{Diese Neigung ist dennoch groß genug, um die Bahn mit hinreichender Sicherheit und Genauigkeit auf der Grundlage von \emph{drei} Beobachtungen zu bestimmen: Tatsächlich kommen die ersten Elemente, die auf diese Weise aus Beobachtungen mit einem Abstand von nur 19 Tagen abgeleitet worden waren (siehe \textsc{von Zach} Monatl. Corresp. Vol. XV. p. 595), bereits denjenigen sehr nahe, die hier aus vier Beobachtungen mit einem Abstand von 162 Tagen abgeleitet werden.}. Wir wählen die folgenden Beobachtungen aus, die in Bremen, Paris, Lilienthal und Mailand von den angesehenen Astronomen \textsc{Olbers}, \textsc{Bouvard}, \textsc{Bessel} und \textsc{Oriani} angestellt wurden:
\begin{center}
\begingroup\fontsize{9}{11}\selectfont\setlength{\tabcolsep}{3pt}
\begin{tabular}{|l|c|l|}
\hline
\shortstack{Mittlere Ortszeit\\der Beobachtung} & Rektaszension & Deklination\\\hline
1807 März 30. $12^{\mathrm h}33^{\mathrm m}17^{\mathrm s}$ & $183^0\,52'\,40''{,}8$ & $11^0\,54'\,27''$ N.\\
\phantom{1807 }Mai 17. $8\quad16\quad5$ & $178\quad36\quad42{,}3$ & $11\quad39\quad46{,}8$ —\\
\phantom{1807 }Juli 11. $10\quad30\quad19$ & $189\quad49\quad7{,}7$ & $3\quad9\quad10{,}1$ —\\
\phantom{1807 }Sept. 8. $7\quad22\quad16$ & $212\quad50\quad3{,}4$ & $8\quad38\quad17{,}0$ S.\\\hline
\end{tabular}\endgroup
\end{center}
Für dieselben Zeitpunkte entnehmen wir den Tafeln der Sonnenbewegungen
\begin{center}
\begingroup\fontsize{9}{11}\selectfont\setlength{\tabcolsep}{2pt}
\begin{tabular}{|l|c|c|c|c|c|}
\hline
& \shortstack{Sonnenlänge vom\\scheinb. Äquin.} & Nutation & \shortstack{Entfernung\\von der Erde} & \shortstack{Breite der\\Sonne} & \shortstack{Scheinbare Schiefe\\der Ekliptik}\\\hline
März 30 & $9^0\,21'\,59''{,}5$ & $+16''{,}8$ & $0{,}999\,6448$ & $+0''{,}23$ & $23^0\,27'\,50''{,}82$\\
Mai 17 & $55\quad56\quad20{,}0$ & $+16{,}2$ & $1{,}0119789$ & $-0{,}63$ & $49{,}83$\\
Juli 11 & $108\quad34\quad53{,}3$ & $+17{,}3$ & $1{,}0165795$ & $-0{,}46$ & $49{,}19$\\
Sept. 8 & $165\quad8\quad57{,}1$ & $+16{,}7$ & $1{,}0067421$ & $+0{,}29$ & $49{,}26$\\\hline
\end{tabular}\endgroup
\end{center}

Die beobachteten Orte des Planeten wurden nun unter Verwendung der scheinbaren Schiefe der Ekliptik in Längen und Breiten umgerechnet, um die Nutation und die Aberration der Fixsterne bereinigt und schließlich nach Abzug der Präzession auf den Beginn des Jahres 1807 bezogen, dann wurden nach den Anweisungen in Art. 72 aus den Orten der Sonne die fiktiven Orte der Erde abgeleitet (um der Parallaxe Rechnung zu tragen) und ihre Längen nach Abzug von Nutation und Präzession auf dieselbe Epoche bezogen; zuletzt wurden die Zeiten vom Jahresbeginn an gezählt und auf den Pariser Meridian bezogen. So ergaben sich die folgenden Zahlen:
\begin{center}
\begingroup\fontsize{9}{11}\selectfont\setlength{\tabcolsep}{2.5pt}
\begin{tabular}{l|c|c|c|c}
$t,\ t',\ t'',\ t'''$ & $89{,}505\,162$ & $137{,}344502$ & $192{,}419502$ & $251{,}288102$\\
$a,\ a',\ a'',\ a'''$ & $178^0\,43'\,38''{,}87$ & $174^0\,1'\,30''{,}08$ & $187^0\,45'\,42''{,}23$ & $213^0\,34'\,15''{,}63$\\
$\beta,\ \beta',\ \beta'',\ \beta'''$ & $12\quad27\quad6{,}16$ & $10\quad8\quad7{,}80$ & $6\quad47\quad25{,}51$ & $4\quad20\quad21{,}63$\\
$l,\ l',\ l'',\ l'''$ & $189\quad21\quad33{,}71$ & $235\quad56\quad0{,}63$ & $288\quad35\quad20{,}32$ & $345\quad9\quad18{,}69$\\
$\mathrm{logg}\,R,R',R'',R'''$ & $9{,}9997990$ & $0{,}005\,1376$ & $0{,}0071739$ & $0{,}0030625$
\end{tabular}\endgroup
\end{center}

\SourceLeaf{238}Daraus leiten wir ab
\[
\begin{array}{ll}
\gamma' = 168^{0}32'41''{,}34, & \delta' = 62^{0}23'4''{,}88 \\
\gamma'' = 173\ \ 5\ 15{,}68, & \delta'' = 100\ 45\ 1{,}40 \\
\log a' = 9{,}952\,6104, & b' = -11{,}009\,449 \\
\log a'' = 9{,}999\,4839, & b'' = -\ 2{,}082\,036
\end{array}
\]
\[
\begin{array}{ll}
\varkappa' = -1{,}083\,306, & \log \lambda = 0{,}072\,8800,\\
\log \mu' = 9{,}713\,9702_n & \\
\varkappa'' = +6{,}322\,006, & \log \lambda''' = 0{,}079\,8512_n,\\
\log \mu'' = 9{,}838\,7061 & 
\end{array}
\]
\[
\begin{array}{ll}
A'D = 37^{0}17'51''{,}50, & A''D = 89^{0}24'11''{,}84,\\
\epsilon = 9^{0}5'5''{,}48 & \\
B'D = -25\ \ 5\ 13{,}38, & B''D = -11\ 20\ 49{,}56.
\end{array}
\]
Nach Abschluss dieser Vorberechnungen gehen wir zur \emph{ersten Hypothese} über. Aus den
Zeitintervallen erhalten wir
\[
\begin{array}{l}
\log k(t'-t) = 9{,}915\,3666 \\
\log k(t''-t') = 9{,}976\,5359 \\
\log k(t'''-t'') = 0{,}005\,4651
\end{array}
\]
und daraus die ersten Näherungswerte
\[
\begin{array}{ll}
\log P' = 0{,}06117, & \log(1+P') = 0{,}33269,\\
\log Q' = 9{,}59087 & \\
\log P'' = 9{,}97107, & \log(1+P'') = 0{,}28681,\\
\log Q'' = 9{,}68097, & 
\end{array}
\]
und hieraus ferner
\[
\begin{array}{ll}
c' = -7{,}68361, & \log d' = 0{,}04666_n \\
c'' = +2{,}20771, & \log d'' = 0{,}12552.
\end{array}
\]
Mit diesen Werten erhält man nach wenigen Versuchen
die folgende Lösung der Gleichungen I. II:
\[
\begin{array}{ll}
x' = 2{,}04856, & z' = 23^{0}38'17'',\\
\log r' = 0{,}34951 & \\
x'' = 1{,}95745, & z'' = 27\ \ 2\ \ 0,\\
\log r'' = 0{,}34194. & 
\end{array}
\]
Aus $z'$, $z''$ und $\epsilon$ erhalten wir $C'C'' = v''-v' = 17^{0}7'5''$: Daraus
sind $v'-v$, $r$, $v'''-v''$, $r'''$ mittels der folgenden Gleichungen zu bestimmen:
\[
\begin{array}{l}
\log r\ \sin(v'-v) = 9{,}74962, \\ \log r\ \sin(v'-v+17^{0}7'5'') = 0{,}07500 \\
\log r'''\sin(v'''-v'') = 9{,}84729, \\ \log r'''\sin(v'''-v'+17^{0}7'5'') = 0{,}10733,
\end{array}
\]
woraus wir erhalten
\SourceLeaf{239}
\[
\begin{array}{ll}
v'-v = 14^{0}14'32'', & \log r = 0{,}35865 \\
v'''-v'' = 18\ 48\ 33, & \log r''' = 0{,}33887.
\end{array}
\]
Schließlich erhält man
\[
\begin{gathered}
\log(\eta 01) = 0{,}00426, \\ \log(\eta 12) = 0{,}00599, \\ \log(\eta 23) = 0{,}00711,
\end{gathered}
\]
und daraus die korrigierten Werte von $P'$, $P''$, $Q'$, $Q''$:
\[
\begin{array}{ll}
\log P' = 0{,}05944, & \log Q' = 9{,}60374 \\
\log P'' = 9{,}97219, & \log Q'' = 9{,}69581,
\end{array}
\]
auf denen die \emph{zweite Hypothese} aufzubauen ist. Ihre wichtigsten Ergebnisse
lauten:
\[
\begin{array}{ll}
c' = -7{,}67820, & \log d' = 0{,}045\,736_n\mathord{\rule[.1ex]{.35em}{.35pt}} \\
c'' = +2{,}21061, & \log d'' = 0{,}126\,054
\end{array}
\]
\[
\begin{array}{ll}
x' = 2{,}03308, & z' = 23^{0}47'54'',\\
\log r' = 0{,}346\,747 & \\
x'' = 1{,}94290, & z'' = 27\ 12\ 25,\\
\log r'' = 0{,}339\,373 & 
\end{array}
\]
\[
C'C'' = v''-v' = 17^{0}8'0''
\]
\[
\begin{array}{ll}
v'-v = 14^{0}21'36'', & \log r = 0{,}354\,687 \\
v'''-v'' = 18\ 50\ 43, & \log r''' = 0{,}334\,564
\end{array}
\]
\[
\begin{gathered}
\log(\eta 01) = 0{,}004\,359, \\ \log(\eta 12) = 0{,}006\,102, \\ \log(\eta 23) = 0{,}007\,280.
\end{gathered}
\]
Hieraus ergeben sich die erneut korrigierten Werte von $P'$, $P''$, $Q'$, $Q''$:
\[
\begin{array}{ll}
\log P' = 0{,}059\,426, & \log Q' = 9{,}604\,749 \\
\log P'' = 9{,}972\,249, & \log Q'' = 9{,}697\,564,
\end{array}
\]
Gehen wir mit diesen Werten zur \emph{dritten Hypothese} über, so ergeben sich
die folgenden Zahlen:
\[
\begin{array}{ll}
c' = -7{,}67815, & \log d' = 0{,}045\,729_n \\
c'' = +2{,}21076, & \log d'' = 0{,}126\,082
\end{array}
\]
\[
\begin{array}{ll}
x' = 2{,}03255, & z' = 23^{0}48'14'',\\
\log r' = 0{,}346\,653 & \\
x'' = 1{,}94235, & z'' = 27\ 12\ 49,\\
\log r'' = 0{,}339\,276 & 
\end{array}
\]
\[
C'C'' = v''-v' = 17^{0}8'4''
\]
\[
\begin{array}{ll}
v'-v = 14^{0}21'49'', & \log r = 0{,}354\,522 \\
v'''-v'' = 18\ 51\ \ 7, & \log r''' = 0{,}334\,290
\end{array}
\]
\[
\begin{gathered}
\log(\eta 01) = 0{,}004\,363, \\ \log(\eta 12) = 0{,}006\,106, \\ \log(\eta 23) = 0{,}007\,290.
\end{gathered}
\]
\SourceLeaf{240}
Berechnet man nun die Entfernungen von der Erde nach den Anweisungen des vorhergehenden Artikels,
so erhält man:
\[
\begin{array}{ll}
\rho' = 1{,}5635, & \rho'' = 2{,}1319 \\
\log \rho \cos\beta = 0{,}09876, & \log \rho'''\cos\beta''' = 0{,}42842 \\
\log \rho \sin\beta = 9{,}44252, & \log \rho'''\sin\beta''' = 9{,}30905 \\
\beta = 12^{0}26'40'', & \beta''' = 4^{0}20'39'' \\
\log \rho = 0{,}10909, & \log \rho''' = 0{,}42967.
\end{array}
\]
Hieraus erhält man
\[
\begin{array}{c|c|c}
 & \text{Zeitkorrekturen} & \text{Korrigierte Zeiten} \\
\hline
\text{I} & 0{,}007\,335 & 89{,}497\,827 \\
\text{II} & 0{,}008\,921 & 137{,}335\,581 \\
\text{III} & 0{,}012\,165 & 192{,}407\,337 \\
\text{IV} & 0{,}015\,346 & 251{,}272\,756
\end{array}
\]
woraus sich die erneut korrigierten Werte der Größen $P'$, $P''$, $Q'$, $Q''$ ergeben
\[
\begin{array}{ll}
\log P' = 0{,}059\,415, & \log Q' = 9{,}604\,782 \\
\log P'' = 9{,}972\,253, & \log Q'' = 9{,}697\,687.
\end{array}
\]
Stellt man schließlich mit diesen neuen Werten die \emph{vierte Hypothese} auf, so ergeben sich
die folgenden Zahlen:
\[
\begin{array}{ll}
c' = -7{,}678\,116, & \log d' = 0{,}045\,723_n \\
c'' = +2{,}210\,773, & \log d'' = 0{,}126\,084
\end{array}
\]
\[
\begin{array}{ll}
x' = 2{,}032\,473, & z' = 23^{0}48'16''{,}7,\\
\log r' = 0{,}346\,638 & \\
x'' = 1{,}942\,281, & z'' = 27\ 12\ 51{,}7,\\
\log r'' = 0{,}339\,263 & 
\end{array}
\]
\[
\begin{gathered}
v''-v' = 17^{0}8'5''{,}1, \\ \tfrac{1}{2}(u''+u') = 176^{0}7'50''{,}5, \\
\tfrac{1}{2}(u''-u') = 4^{0}33'23''{,}6
\end{gathered}
\]
\[
\begin{array}{ll}
v'-v = 14^{0}21'51''{,}9, & \log r = 0{,}354\,503 \\
v'''-v'' = 18\ 51\ \ 9{,}5, & \log r''' = 0{,}334\,263.
\end{array}
\]
Diese Zahlen weichen so wenig von denen der dritten Hypothese ab,
dass wir nun mit Sicherheit zur Bestimmung der Elemente selbst übergehen
können. Zuerst bestimmen wir die Lage der Bahnebene. Nach den Anweisungen in Art.\ 149
erhält man aus $\gamma'$, $u'$ und $A'C' = \delta'-z'$ die Bahnneigung
$= 7^{0}8'14''{,}8$, die Länge des aufsteigenden Knotens $103^{0}16'37''{,}2$, das Breitenargument
der zweiten Beobachtung $94^{0}36'4''{,}9$, \SourceLeaf{241}also die Länge in der Bahn $197^{0}52'42''{,}1$; ebenso erhält man aus $\gamma''$, $u''$
und $A''C'' = \delta''-z''$ die Bahnneigung $= 7^{0}8'14''{,}8$,
die Länge des aufsteigenden Knotens $103^{0}16'37''{,}5$, das Breitenargument
der dritten Beobachtung $111^{0}44'9''{,}7$, also die Länge in der Bahn
$215^{0}0'47''{,}2$. Damit beträgt die Länge in der Bahn für die erste Beobachtung
$183^{0}30'50''{,}2$, für die vierte $233^{0}51'56''{,}7$. Bestimmt man nun aus $t'''-t$,
$r$, $r'''$ und $v'''-v = 50^{0}21'6''{,}5$ die Abmessungen der Bahn,
so ergibt sich:
\begin{center}
\begingroup\fontsize{9}{11}\selectfont\setlength{\tabcolsep}{3pt}
\begin{tabular}{@{}lr@{}}
Wahre Anomalie für den ersten Ort . . . . . . . . . . . . . . . . . & $293^{0}33'43''{,}7$ \\
Wahre Anomalie für den vierten Ort . . . . . . . . . . . . . . . . . & $343\ \ 54\ \ 50{,}2$ \\
Daraus die Länge des Perihels . . . . . . . . . . . . . . . . . . . . . . . & $249\ \ 57\ \ \ 6{,}5$ \\
Mittlere Anomalie für den ersten Ort . . . . . . . . . . . . . . . . & $302\ \ 33\ \ 32{,}6$ \\
Mittlere Anomalie für den vierten Ort . . . . . . . . . . . . . . . & $346\ \ 32\ \ 25{,}2$ \\
Mittlere siderische Tagesbewegung . . . . . . . . . . . . . . . . . & $978''{,}7216$ \\
Mittlere Anomalie für den Beginn des Jahres 1807 . . . . . . . . . . . & $278^{0}13'39''{,}1$ \\
Mittlere Länge für dieselbe Epoche . . . . . . . . . . . . & $168\ \ 10\ \ 45{,}6$ \\
Winkel $\varphi$ . . . . . . . . . . . . . . . . . . . . . . . . . . . . . . . . . . & $5\ \ \ 2\ \ 58{,}1$ \\
Logarithmus der großen Halbachse . . . . . . . . . . . . . . . . . . & $0{,}372\,898.$
\end{tabular}\endgroup
\end{center}

Berechnet man mit diesen Elementen für die korrigierten Zeiten $t$, $t'$, $t''$, $t'''$
die geozentrischen Orte des Planeten, so stimmen die vier Längen mit $a$, $a'$,
$a''$, $a'''$ und die beiden mittleren Breiten mit $\beta'$, $\beta''$ bis auf
eine Zehntelbogensekunde überein; für die erste und die letzte Breite ergeben sich dagegen
$12^{0}26'43''{,}7$ und $4^{0}20'40''{,}1$, wobei die erste um $22''{,}4$
zu klein, die letzte um $18''{,}5$ zu groß ist. Lässt man jedoch die übrigen Elemente unverändert
und vergrößert lediglich die Bahnneigung um $6''$ und verkleinert die Knotenlänge um $4' 40''$,
so verteilen sich die Fehler auf alle Breiten und sinken auf wenige
Bogensekunden, während die Längen nur sehr geringe Fehler aufweisen,
die ebenfalls fast ganz verschwinden, wenn man darüber hinaus
die Epoche der Länge um $2''$ verkleinert.
\par\medskip\begin{center}\rule{.24\linewidth}{.35pt}\end{center}
